\documentclass[acmsmall,nonacm]{acmart}

%% ── packages ──────────────────────────────────────────────────────
\usepackage{booktabs}
\usepackage{listings}
\usepackage{xcolor}
\usepackage{amsmath}
\usepackage{amssymb}
\usepackage{hyperref}
\usepackage{subcaption}
\usepackage{algorithm}
\usepackage{algorithmic}

%% ── listings style ────────────────────────────────────────────────
\definecolor{codebg}{gray}{0.96}
\definecolor{codegreen}{rgb}{0.13,0.55,0.13}
\definecolor{codegray}{rgb}{0.5,0.5,0.5}
\lstset{
  basicstyle=\ttfamily\small,
  backgroundcolor=\color{codebg},
  keywordstyle=\color{blue},
  commentstyle=\color{codegreen},
  stringstyle=\color{red!70!black},
  numbers=left,
  numberstyle=\tiny\color{codegray},
  numbersep=5pt,
  breaklines=true,
  frame=single,
  framesep=4pt,
  xleftmargin=8pt,
  xrightmargin=4pt,
  columns=fullflexible,
  keepspaces=true,
  showstringspaces=false,
  literate={→}{$\rightarrow$}1 {×}{$\times$}1 {≈}{$\approx$}1,
}
\lstdefinelanguage{C89}{
  keywords={int, double, void, static, const, struct, typedef, enum,
            if, else, for, while, return, switch, case, break, default,
            sizeof, malloc, calloc, free, memcpy},
  morecomment=[l]{//},
  morecomment=[s]{/*}{*/},
  morestring=[b]",
}

%% ── metadata ──────────────────────────────────────────────────────
\title{Tag-System Scheduled Transformer Inference:\\
  A 2-Tag System as the Control Model for LLM Generation}

\author{Alvaro Videla}
\affiliation{%
  \institution{Microsoft}
}

\begin{abstract}
We present a transformer inference engine whose control flow is
governed entirely by a 2-tag system in the sense of Post~(1943).
The engine maintains a two-symbol word $[\mathit{OP}, \mathit{STATE}]$
and advances by reading the operation symbol, executing the
corresponding production rule (one neural-network sub-operation),
and emitting the next $[\mathit{OP}', \mathit{STATE}']$ pair.
For a 24-layer Qwen~2.5 0.5B model, each generated token requires
exactly 315~tag steps.

The production rules encode sign-plus-magnitude matrix--vector
products in the log domain, where weights are stored as
$(\ln|w|, \operatorname{sign}(w))$ pairs and accumulation uses the
log-sum-exp identity.  We apply five optimizations drawn from
the classical computer science literature:
(a)~a peephole-rewritten fast $\exp$ approximation (Aho et~al., Dragon Book),
(b)~semigroup tree reduction for instruction-level parallelism (Stepanov \& McJones),
(c)~constraint-propagation pruning at threshold $e^{-20}$ (Dechter),
(d)~Grand Central Dispatch parallel column decomposition, and
(e)~a Metal GPU backend with a SIMD-cooperative kernel using
half-precision magnitudes and packed sign bits.

We further introduce the \texttt{.tag} model format: a GPU-native,
\texttt{mmap}-friendly binary encoding with page-aligned f16
magnitudes and packed sign bits, enabling zero-copy Metal buffer
registration via \texttt{newBufferWithBytesNoCopy}.  The format
reduces model size from 4.8\,GB (f64 \texttt{.eml}) to 2.0\,GB
(58\% reduction) and eliminates the GPU upload phase entirely.

On an Apple M4~Max, the engine generates correct English text at
${\sim}12$~tok/s wall-clock, with the matmul kernel achieving a
$97.8\times$ speedup over the baseline CPU implementation.
The entire system---tag-system scheduler, five optimization
layers, GPU bridge, and \texttt{mmap} loader---is implemented in
approximately 1700~lines of~C and 350~lines of Objective-C.

\end{abstract}

\keywords{tag system, Post production system, transformer inference,
  sign-magnitude arithmetic, Metal GPU, mmap, zero-copy}

\begin{document}
\maketitle

%% ═════════════════════════════════════════════════════════════════
\section{Introduction}
\label{sec:intro}

Post's tag systems~\cite{post1943} are among the simplest known
Turing-complete computational models.  A $m$-tag system operates on a
finite alphabet $A$ with a production function $P : A \to A^*$: at
each step, the machine reads the leftmost symbol~$a$ of the word,
appends $P(a)$ to the right end, and deletes the leftmost~$m$
symbols.  Cocke and Minsky~\cite{cocke1964} proved that 2-tag systems
($m=2$) are universal.

Modern large language models (LLMs) are typically implemented as
imperative programs with nested loops over layers, attention heads,
and sequence positions.  The control flow is implicit in the loop
structure.  We ask: \emph{can a tag system serve as the explicit
scheduling mechanism for transformer inference?}

We answer affirmatively by constructing \textsc{emilio-tag}, an
inference engine for Qwen~2.5 0.5B~Instruct~\cite{qwen2025} in
which:

\begin{enumerate}
\item The computation state is a two-symbol word
  $[\mathit{OP}, \mathit{STATE}]$, where $\mathit{OP}$ determines
  which neural-network sub-operation to execute and $\mathit{STATE}$
  carries the mutable activation vectors.

\item Each of 16 production rules performs one transformer
  sub-operation (embedding lookup, RMSNorm, QKV projection,
  RoPE, attention, FFN, etc.) and emits the next operation symbol.

\item The word length is invariant at~2.  The tag system halts
  after exactly $1 + 24 \times 13 + 2 = 315$ steps per token.

\item The production rules are \emph{parameterised} by model
  weights: $P(\textsc{qkv\_matmul}, \text{layer}=k)$ uses the
  QKV weight matrix of layer~$k$.
\end{enumerate}

The tag system governs sequential control flow---which operation
runs after which---while the evaluation of each production rule
uses conventional CPU or GPU arithmetic.  We call this architecture
\emph{tag-system-scheduled, natively-evaluated}.

The rest of this paper describes the formal model
(Section~\ref{sec:model}), the sign-magnitude matmul kernel and its
five optimization layers (Section~\ref{sec:optimizations}), the Metal
GPU backend (Section~\ref{sec:gpu}), the \texttt{.tag} model format
with \texttt{mmap} and zero-copy GPU registration
(Section~\ref{sec:format}), experimental results
(Section~\ref{sec:results}), and related work
(Section~\ref{sec:related}).


%% ═════════════════════════════════════════════════════════════════
\section{The 2-Tag System Model}
\label{sec:model}

\subsection{Formal Definition}

A 2-tag system is a triple $(m, A, P)$ where:
\begin{itemize}
\item $m = 2$ is the \emph{deletion number}.
\item $A$ is a finite \emph{alphabet}.
\item $P : A \to A^*$ maps each symbol to a \emph{production}
  (a finite string over~$A$).
\end{itemize}

The machine maintains a \emph{word} $w \in A^*$.  At each step:
\begin{enumerate}
\item Read the leftmost symbol $a = w_0$.
\item If $a$ is a halting symbol, stop.
\item Append $P(a)$ to the right end of~$w$.
\item Delete the leftmost $m = 2$ symbols.
\end{enumerate}

\subsection{Instantiation for Transformer Inference}

We define the alphabet for a transformer with $L$~layers:
\begin{align*}
A = \{\, &\textsc{embed},\; \textsc{rmsnorm\_attn},\;
  \textsc{qkv\_matmul},\; \textsc{bias\_add},\\
  &\textsc{rope},\; \textsc{kv\_store},\;
  \textsc{attention},\; \textsc{out\_proj},\\
  &\textsc{residual\_attn},\; \textsc{ffn\_norm},\;
  \textsc{gate\_up},\; \textsc{silu\_mul},\\
  &\textsc{down\_proj},\; \textsc{residual\_ffn},\;
  \textsc{final\_norm},\; \textsc{lm\_head},\\
  &\textsc{halt},\; \textsc{state}\,\}
\end{align*}

Each symbol carries a layer index $\ell \in \{0, \dots, L-1\}$.
The word is always exactly two symbols:
\[
  w = [\mathit{OP}_\ell,\; \textsc{state}]
\]

The \textsc{state} symbol carries a mutable payload: the current
activation vector~$\mathbf{x}$, a residual vector, the sequence
position, and the current layer index.

\subsection{Production Rules}

Each production $P(\mathit{OP}_\ell)$ performs one sub-operation
and determines the next operation symbol.  The key rules are:

\medskip
\noindent\textbf{Embedding.}
$P(\textsc{embed})$: look up $\mathbf{x} \leftarrow E[\text{token\_id}]$,
emit $[\textsc{rmsnorm\_attn}_0, \textsc{state}]$.

\medskip
\noindent\textbf{Layer transition.}
$P(\textsc{residual\_ffn}_\ell)$: compute
$\mathbf{x} \leftarrow \mathbf{x}_\text{residual} + \mathbf{x}_\text{ffn}$.
If $\ell < L-1$, emit $[\textsc{rmsnorm\_attn}_{\ell+1}, \textsc{state}]$;
otherwise emit $[\textsc{final\_norm}, \textsc{state}]$.

\medskip
\noindent\textbf{Halting.}
$P(\textsc{lm\_head})$: compute output logits via matmul, set
$\text{argmax}$, emit $[\textsc{halt}, \textsc{state}]$.

\subsection{Step Count Analysis}

For one token through a 24-layer model:
\begin{align*}
\text{steps} &= 1_{\text{embed}} + L \times 13_{\text{layer ops}}
  + 1_{\text{final\_norm}} + 1_{\text{lm\_head}} \\
  &= 1 + 24 \times 13 + 2 = 315
\end{align*}

The 13 layer operations are: \textsc{rmsnorm\_attn},
\textsc{qkv\_matmul}, \textsc{bias\_add}, \textsc{rope},
\textsc{kv\_store}, \textsc{attention}, \textsc{out\_proj},
\textsc{residual\_attn}, \textsc{ffn\_norm}, \textsc{gate\_up},
\textsc{silu\_mul}, \textsc{down\_proj}, \textsc{residual\_ffn}.


%% ═════════════════════════════════════════════════════════════════
\section{Sign-Magnitude Arithmetic and Optimizations}
\label{sec:optimizations}

\subsection{Log-Domain Sign-Magnitude Representation}

Weights are stored as pairs $(\ln|w_k|, \operatorname{sign}(w_k))$.
A matrix--vector product $\mathbf{y} = W\mathbf{x}$ becomes:

\begin{equation}
y_j = \sum_{k} \operatorname{sign}(w_{jk}) \cdot
  \operatorname{sign}(x_k) \cdot
  \exp\!\bigl(\ln|w_{jk}| + \ln|x_k|\bigr)
\label{eq:smmatmul}
\end{equation}

At runtime, the activation $\mathbf{x}$ is converted to
$(\ln|x_k|, \operatorname{sign}(x_k))$ before each matmul.  The
inner products are computed via log-sum-exp with a running maximum
for numerical stability.

\subsection{Optimization 1: Fast Exponential (Dragon Book)}

The inner loop calls $\exp()$ approximately $K$ times per dot
product, where $K$ is the inner dimension (896 for most layers,
4864 for \textsc{down\_proj}).  We replace the libm \texttt{exp()}
with a range-reduced 5th-degree polynomial~\cite{dragonbook}:

\begin{enumerate}
\item \textbf{Range reduction}: $x = n \cdot \ln 2 + r$,
  where $n = \lfloor x / \ln 2 \rfloor$ and $0 \le r < \ln 2$.
\item \textbf{Polynomial}: $\exp(r) \approx 1 + r(1 + r(0.5 +
  r(\tfrac{1}{6} + r(\tfrac{1}{24} + \tfrac{r}{120}))))$.
\item \textbf{Reconstruction}: $2^n$ via IEEE~754 exponent bit
  manipulation (zero floating-point operations).
\end{enumerate}

The relative error is below $2 \times 10^{-7}$, well within the
tolerance of the sign-magnitude representation.  This is a
\emph{peephole optimization} in the Dragon Book sense: a local
pattern (expensive \texttt{exp} call) is replaced by a cheaper
equivalent without changing the surrounding code.

\subsection{Optimization 2: Pruning Threshold (Dechter)}

Terms where $\ln|w_{jk}| + \ln|x_k| - \max < -20$ contribute at most
$e^{-20} \approx 2 \times 10^{-9}$ to the sum.  Over $K = 896$
terms, the maximum accumulated error is $K \cdot e^{-20} \approx
1.8 \times 10^{-6}$.  We skip these terms entirely.

This is constraint propagation in the sense of
Dechter~\cite{dechter}: propagating the constraint ``this term's
contribution is negligible'' to prune the computation.  The
threshold was tightened from $-30$ to $-20$ after analyzing the
error budget.

\subsection{Optimization 3: Semigroup Tree Reduction (Stepanov)}

The accumulation $\texttt{acc} = \texttt{acc} + f(x_k)$ is a linear
dependency chain of depth~$K$.  Since addition is associative (it
forms a semigroup~\cite{stepanov2009}), we can reparenthesise:

\begin{lstlisting}[language=C89,caption={4-way ILP accumulation}]
double a0 = 0, a1 = 0, a2 = 0, a3 = 0;
for (k = 0; k < K; k += 4) {
    a0 += fast_exp(s[k])   * sign[k];
    a1 += fast_exp(s[k+1]) * sign[k+1];
    a2 += fast_exp(s[k+2]) * sign[k+2];
    a3 += fast_exp(s[k+3]) * sign[k+3];
}
acc = (a0 + a1) + (a2 + a3);  // tree reduce
\end{lstlisting}

The final tree-reduce \texttt{(a0+a1)+(a2+a3)} has depth 2 instead
of~$K$.  The same 4-way decomposition is applied to the max-finding
pass.

\subsection{Optimization 4: GCD Parallel Dispatch}

The column loop of the matrix--vector product is embarrassingly
parallel: each output element $y_j$ depends on a different column
of~$W$.  We split the columns across $N$ worker threads using macOS
Grand Central Dispatch (\texttt{dispatch\_apply\_f}):

\begin{lstlisting}[language=C89,caption={GCD parallel dispatch}]
dispatch_apply_f(nchunks,
    dispatch_get_global_queue(
        QOS_CLASS_USER_INTERACTIVE, 0),
    &ctx, matmul_dispatch_work);
\end{lstlisting}

Thread count is auto-detected via \texttt{sysconf(\_SC\_NPROCESSORS\_ONLN)}
and overridable with \texttt{EML\_THREADS}.  The GCD thread pool
has zero creation/destruction overhead across the ${\sim}3000$
matmul calls per 8-token generation.


%% ═════════════════════════════════════════════════════════════════
\section{Metal GPU Backend}
\label{sec:gpu}

\subsection{Architecture}

The tag system determines \emph{which} operation to execute; the GPU
determines \emph{how} the heaviest operations (matmuls) are evaluated.
When a production rule calls \texttt{build\_matmul\_sm} and a GPU
weight buffer exists (\texttt{gpu\_weight\_id} $\ge 0$), the call is
dispatched to the Metal compute shader.  All other production rules
(RMSNorm, softmax, RoPE, residual adds) run on the CPU.

This separation is principled: the tag system is the \emph{control
model} and the GPU is the \emph{evaluation model}.  The tag system
still executes all 315 sequential steps per token---parallelism
exists only \emph{within} a single production rule's evaluation.

\subsection{SIMD-Cooperative Kernel}

The Metal shader assigns one SIMD group (32~threads) per output
element $(i, j)$.  Each thread processes $\lceil K/32 \rceil$
elements of the inner dimension and contributes to a
\texttt{simd\_sum} reduction:

\begin{lstlisting}[language=C89,caption={Metal kernel core (simplified)}]
// 4-way unrolled inner loop
for (; k + 96 < inner; k += 128) {
    float m0 = ln_a_mag[ak]   + half(ln_b_mag[bk]);
    float m1 = ln_a_mag[ak+32]+ half(ln_b_mag[bk+32]);
    float m2 = ln_a_mag[ak+64]+ half(ln_b_mag[bk+64]);
    float m3 = ln_a_mag[ak+96]+ half(ln_b_mag[bk+96]);
    // sign extraction from packed u32
    float sb0 = 1-2*float((w0>>lane)&1);
    acc0 += exp(m0) * (sign_a[ak] * sb0);
    ...
}
float total = simd_sum(acc0+acc1+acc2+acc3);
\end{lstlisting}

\subsection{Weight Storage on GPU}

Weight magnitudes are stored as f16 (half-precision), and signs
are packed as single bits into u32 words (bit~$k$ of word~$\lfloor k/32 \rfloor$
is 1 iff $w_k < 0$).  This achieves a 23\% bandwidth reduction
over storing magnitudes as f32 and signs as f32.

The activation vector~$\mathbf{x}$ is converted to
$(\ln|x_k|, \operatorname{sign}(x_k))$ on the CPU and written to
shared Metal buffers.  The result is read back as f32 and
converted to f64 for the next tag step.


%% ═════════════════════════════════════════════════════════════════
\section{The \texttt{.tag} Model Format}
\label{sec:format}

\subsection{Motivation}

The original \texttt{.eml} format stores weights as f64
sign-magnitude pairs.  Loading the 0.5B model requires:
\begin{enumerate}
\item Reading 4.8\,GB of f64 data from disk (\(\sim\)10\,s).
\item Converting each weight from f64 to f16 magnitude and
  packing signs into u32 words.
\item Copying the converted data to Metal GPU buffers
  (\(\sim\)6\,s).
\end{enumerate}

We designed the \texttt{.tag} format to eliminate steps 2 and~3.

\subsection{Format Layout}

The \texttt{.tag} file has four sections:

\begin{description}
\item[Header (84 bytes):] Magic \texttt{TAG1}, version~(u32),
  model configuration (10$\times$u32 + 2$\times$f64),
  tensor count~(u32), and offsets to the tokenizer, tensor
  directory, and data sections.

\item[Tokenizer:] Vocabulary strings and BPE merge table, stored
  as length-prefixed byte sequences.

\item[Dense arrays:] Token embeddings, layer norms, and biases
  as f64 arrays with u64 length prefixes.

\item[Tensor data (page-aligned):] f16 magnitudes and packed u32
  sign bits for all 97~weight matrices, laid out contiguously
  starting at a 16\,KB-aligned offset.
\end{description}

A \emph{tensor directory} between the dense arrays and tensor data
provides the inner dimension, column count, magnitude offset, and
sign offset for each tensor.

\subsection{Zero-Copy GPU Registration}

The tensor data section is page-aligned (16\,KB on Apple Silicon)
so that the \texttt{mmap}'d region can be wrapped directly as a
Metal buffer:

\begin{lstlisting}[language=C89,caption={Zero-copy Metal buffer from mmap}]
id<MTLBuffer> buf = [device
    newBufferWithBytesNoCopy:(void *)ptr
        length:nbytes
        options:MTLResourceStorageModeShared
        deallocator:nil];
\end{lstlisting}

On Apple Silicon's unified memory architecture, this avoids any data
copy: the GPU reads directly from the \texttt{mmap}'d virtual pages.
The \texttt{nil} deallocator indicates that the mmap owns the
underlying memory.  If the pointer is not page-aligned (e.g., on
non-Apple hardware), the implementation falls back to a copy.

\subsection{Loading Procedure}

\begin{enumerate}
\item \texttt{mmap()} the entire \texttt{.tag} file with
  \texttt{MAP\_PRIVATE} and \texttt{PROT\_READ}.
\item Parse the header and tokenizer from the mapped region
  (pointer arithmetic, no I/O).
\item Copy dense f64 arrays (embeddings, norms, biases) to
  heap-allocated buffers.
\item Iterate the tensor directory: for each of the 97~tensors,
  call \texttt{metal\_register\_weights\_nocopy()} with pointers
  into the mapped data section.
\end{enumerate}

The kernel does not touch the tensor data during loading---it
is faulted in by the GPU on first access.

\subsection{File Size Comparison}

\begin{table}[h]
\centering
\caption{Model format comparison for Qwen 2.5 0.5B Instruct.}
\label{tab:format}
\begin{tabular}{lrrl}
\toprule
\textbf{Format} & \textbf{Size} & \textbf{Precision} & \textbf{Loading} \\
\midrule
\texttt{.gguf} (q8) & 0.5\,GB & int8 & standard \\
\texttt{.eml} (v2)  & 4.8\,GB & f64  & \texttt{fread} + convert \\
\texttt{.tag} (v1)  & 2.0\,GB & f16+bits & \texttt{mmap} zero-copy \\
\bottomrule
\end{tabular}
\end{table}


%% ═════════════════════════════════════════════════════════════════
\section{Experimental Results}
\label{sec:results}

All experiments use Qwen~2.5 0.5B Instruct (24~layers,
$d_\text{model}=896$, $d_\text{ff}=4864$, vocab$=151936$) on an
Apple M4~Max.  The prompt is ``What is a tag system?'' with 8
generated tokens.  Output in all configurations: ``A tag system
is a type of data''.

\subsection{Optimization Progression}

\begin{table}[h]
\centering
\caption{Cumulative effect of optimizations on matmul time and wall
  clock (8 tokens generated).}
\label{tab:optim}
\begin{tabular}{lrrr}
\toprule
\textbf{Configuration} & \textbf{Matmul} & \textbf{Wall} & \textbf{Speedup} \\
\midrule
Baseline (na\"ive C)      & 42,858\,ms & 63\,s   & $1.0\times$ \\
+ fast\_exp + ILP + prune  & 29,774\,ms & 43\,s   & $1.4\times$ \\
+ GCD 16 threads           &  3,611\,ms & 16\,s   & $11.9\times$ \\
+ Metal GPU                &    438\,ms & 16.5\,s & $97.8\times$ \\
+ \texttt{.tag} mmap       &    668\,ms & 12.9\,s & --- \\
\bottomrule
\end{tabular}
\end{table}

The matmul speedup from baseline to GPU is $97.8\times$.  Note that
the \texttt{.tag} format improves wall-clock time (12.9\,s vs.\ 16.5\,s)
by eliminating the f64$\to$f16 conversion and GPU upload phase, while
inference time remains comparable.

\subsection{Time Breakdown by Production Rule}

Table~\ref{tab:profile} shows the time spent in each production rule
when using the Metal GPU backend with \texttt{.tag} loading.

\begin{table}[h]
\centering
\caption{Production rule profile (Metal GPU, \texttt{.tag} format,
  8 generated tokens, 32 forward passes).}
\label{tab:profile}
\begin{tabular}{lrrr}
\toprule
\textbf{Rule} & \textbf{Count} & \textbf{Time (ms)} & \textbf{\%} \\
\midrule
\textsc{embed}          &  32 &   0.0 &  0.0 \\
\textsc{rmsnorm\_attn}  & 768 &   1.1 &  0.2 \\
\textsc{qkv\_matmul}    & 768 & 111.4 & 16.7 \\
\textsc{bias\_add}      & 768 &   0.4 &  0.1 \\
\textsc{rope}           & 768 &   0.2 &  0.0 \\
\textsc{kv\_store}      & 768 &   0.4 &  0.1 \\
\textsc{attention}      & 768 &  11.1 &  1.7 \\
\textsc{out\_proj}      & 768 & 103.5 & 15.5 \\
\textsc{residual\_attn} & 768 &   0.2 &  0.0 \\
\textsc{ffn\_norm}      & 768 &   0.9 &  0.1 \\
\textsc{gate\_up}       & 768 & 198.6 & 29.8 \\
\textsc{silu\_mul}      & 768 &   5.5 &  0.8 \\
\textsc{down\_proj}     & 768 & 159.8 & 23.9 \\
\textsc{residual\_ffn}  & 768 &   0.2 &  0.0 \\
\textsc{final\_norm}    &  32 &   0.0 &  0.0 \\
\textsc{lm\_head}       &  32 &  74.4 & 11.1 \\
\midrule
\textbf{Total}          &     & 667.5 & 100.0 \\
\bottomrule
\end{tabular}
\end{table}

The four matmul-heavy production rules (\textsc{gate\_up},
\textsc{down\_proj}, \textsc{qkv\_matmul}, \textsc{out\_proj})
account for 85.9\% of compute time.  Non-matmul operations
(norms, residual adds, softmax, RoPE) are negligible at 2.2\%.

\subsection{Format Loading Comparison}

\begin{table}[h]
\centering
\caption{Loading performance: \texttt{.eml} vs.\ \texttt{.tag}.}
\label{tab:loading}
\begin{tabular}{lrrr}
\toprule
\textbf{Format} & \textbf{File} & \textbf{Wall} & \textbf{CPU time} \\
\midrule
\texttt{.eml} (fread + upload) & 4.8\,GB & 16.1\,s & 13.9\,s \\
\texttt{.tag} (mmap zero-copy) & 2.0\,GB & 12.9\,s & 11.2\,s \\
\bottomrule
\end{tabular}
\end{table}


%% ═════════════════════════════════════════════════════════════════
\section{Discussion}
\label{sec:discussion}

\subsection{Tag System as Control Model}

The tag system is the \emph{scheduling layer}: it determines the
order of operations.  The 315 sequential steps per token are a
faithful representation of the transformer's dataflow dependencies.
No step can be reordered without violating data dependencies---the
tag system encodes the \emph{minimum} sequential structure.

Within each production rule, however, the evaluation is conventional:
SIMD instructions, GPU warps, and multi-core parallelism.  We call
this \emph{tag-system-scheduled, natively-evaluated}: the control
model is Post's (1943) 2-tag system, the evaluation model is
contemporary hardware.

\subsection{Why a Tag System?}

Beyond theoretical interest, the tag system provides:
\begin{itemize}
\item \textbf{Explicit dataflow.}  Every operation and its
  dependencies are visible in the production rules.  There are no
  hidden control paths.
\item \textbf{Deterministic step count.}  Exactly 315 steps per
  token, invariant of input.  This makes performance analysis
  trivial: profile each production rule independently.
\item \textbf{Composability.}  New operations (e.g., mixture of
  experts, sliding-window attention) are added by defining new
  production rules and extending the alphabet.
\item \textbf{Formal foundation.}  The system's computational
  universality is guaranteed by Cocke and Minsky's
  theorem~\cite{cocke1964}.
\end{itemize}

\subsection{Limitations}

The sign-magnitude representation in f64 is space-inefficient
compared to standard quantization formats (q4, q8).  The
\texttt{.tag} format mitigates this with f16 storage, but a 0.5B
model still occupies 2.0\,GB.  The CPU path remains necessary for
non-matmul operations, creating a CPU--GPU synchronization boundary
at every production rule that involves a matmul.


%% ═════════════════════════════════════════════════════════════════
\section{Related Work}
\label{sec:related}

\paragraph{Tag systems and universality.}
Post~\cite{post1943} introduced tag systems as a model of
computation.  Cocke and Minsky~\cite{cocke1964} proved 2-tag
universality.  De~Mol~\cite{demol2008} surveys the historical
development.  To our knowledge, no prior work uses tag systems
for scheduling neural network computations.

\paragraph{EML operator.}
Odrzywołek~\cite{odrzywolek2026} showed that
$\mathrm{eml}(x,y) = e^x - \ln y$ can derive all elementary
functions.  Our sign-magnitude encoding is related: the inner
product is computed as $\exp(\ln|a| + \ln|b|) \cdot \operatorname{sign}(a)
\cdot \operatorname{sign}(b)$, which factors through $\exp$ and $\ln$.

\paragraph{Efficient inference.}
\texttt{llama.cpp}~\cite{llamacpp} and \texttt{ggml} provide
highly optimized quantized inference.  Our work is not competitive
on throughput; it explores a different point in the design space
where the control model is an explicit formal system.

\paragraph{Memory-mapped model loading.}
The \texttt{gguf} format supports \texttt{mmap}-based loading.
\texttt{safetensors}~\cite{safetensors} also uses memory mapping.
Our \texttt{.tag} format extends this to zero-copy GPU buffer
registration via Metal's \texttt{newBufferWithBytesNoCopy}, avoiding
even the kernel-to-GPU copy on unified memory architectures.

\paragraph{Classical optimization sources.}
Our optimizations draw explicitly from:
Aho et~al.~\cite{dragonbook} (peephole optimization),
Stepanov and McJones~\cite{stepanov2009} (semigroup reduction),
and Dechter~\cite{dechter} (constraint propagation for pruning).


%% ═════════════════════════════════════════════════════════════════
\section{Conclusion}
\label{sec:conclusion}

We have shown that a 2-tag system can serve as the explicit control
model for transformer inference.  The engine generates correct text
from a 0.5B-parameter model using 315 tag steps per token, with
each step corresponding to one transformer sub-operation.  Five
classical optimization techniques---peephole rewriting, semigroup
reduction, constraint-propagation pruning, multi-core dispatch,
and GPU offloading---bring the matmul kernel from 42.9\,s to
0.44\,s ($97.8\times$).  The \texttt{.tag} model format with
\texttt{mmap} and zero-copy Metal registration reduces file size
by 58\% and loading time by 20\%.

The system demonstrates that even the simplest Turing-complete
formalism---a 2-tag system with deletion number $m=2$ and 18
symbols---can organize the control flow of a modern neural network.
Speed is not the primary contribution; rather, the work establishes
that formal models of computation from the 1940s remain applicable
as scheduling substrates for contemporary AI systems.


%% ═════════════════════════════════════════════════════════════════
\bibliographystyle{ACM-Reference-Format}

\begin{thebibliography}{10}

\bibitem{post1943}
E.~L. Post.
\newblock Formal reductions of the general combinatorial decision problem.
\newblock {\em American Journal of Mathematics}, 65(2):197--215, 1943.

\bibitem{cocke1964}
J.~Cocke and M.~Minsky.
\newblock Universality of tag systems with {P}=2.
\newblock {\em Journal of the ACM}, 11(1):15--20, 1964.

\bibitem{qwen2025}
Qwen Team.
\newblock Qwen2.5 Technical Report.
\newblock \emph{arXiv preprint arXiv:2412.15115}, 2024.

\bibitem{odrzywolek2026}
A.~Odrzywołek.
\newblock All elementary functions from a single binary operator.
\newblock arXiv:2603.21852, 2026.

\bibitem{dragonbook}
A.~V. Aho, M.~S. Lam, R.~Sethi, and J.~D. Ullman.
\newblock {\em Compilers: Principles, Techniques, and Tools}.
\newblock Pearson, 2nd edition, 2006.

\bibitem{stepanov2009}
A.~A. Stepanov and P.~McJones.
\newblock {\em Elements of Programming}.
\newblock Addison-Wesley, 2009.

\bibitem{dechter}
R.~Dechter.
\newblock {\em Constraint Processing}.
\newblock Morgan Kaufmann, 2003.

\bibitem{demol2008}
L.~De~Mol.
\newblock Tag systems and {C}ollatz-like functions.
\newblock {\em Theoretical Computer Science}, 390(1):92--101, 2008.

\bibitem{llamacpp}
G.~Gerganov et~al.
\newblock llama.cpp: {LLM} inference in {C/C++}.
\newblock \url{https://github.com/ggerganov/llama.cpp}, 2023--2026.

\bibitem{safetensors}
Hugging Face.
\newblock Safetensors: Simple, safe way to store and distribute tensors.
\newblock \url{https://github.com/huggingface/safetensors}, 2023.

\end{thebibliography}

\end{document}
